\sporkSection{Appendix: Ping-Pong Kernel}{sec:PingPongAppendix}

We constrast the \texttt{sm\_90a} cooperative GEMM kernel with the \texttt{sm\_90a} ping-pong GEMM kernel
(Figure~\ref{fig:PingPongListings}).
The cooperative kernel names 2 ``consumer'' warpgroups (8 warps total)
per CTA to compute an output \texttt{C} tile,
and the ping-pong kernel subdivides the threads further into 2 consumer warpgroups,
each computing their own output tile.
In both cases, the consumers alternate between a main loop (accumulation) and an epilogue
(stage output tile to \texttt{C\_smem}, then store to \texttt{C}).
The ping-pong kernel has the advantage that one consumer warpgroup is in the main loop
while the other is in the epilogue, hiding the memory-bound epilogue's latency.
Although the cooperative kernel was fairly straightforward,
the ping-pong kernel exposed some limitations of \sporkName{}.

% Both listings side by side at one shared small font size; the narrower column is coop.
\begin{figure}[tbp]
\sporkCodeRule
\fontsize{8}{9.6}\selectfont
\begin{minipage}[t]{0.465\linewidth}
\vspace{0pt}
\sporkListing{coop}
\end{minipage}%
\hfill
\begin{minipage}[t]{0.515\linewidth}
\vspace{0pt}
\sporkListing{ping_pong}
\end{minipage}
\caption{Cooperative (left) and ping-pong (right) \texttt{sm\_90a} GEMM pseudocode.
We use real \sporkName{} syntax, but use scalars to stand-in for tiles and omit variable declarations
and detailed \texttt{wgmma} and TMA synchronization.}
\label{fig:PingPongListings}
\sporkCodeRule
\end{figure}


\sporkMainKey{Ping-pong Loop Unrolling:}
The \sporkRed{\texttt{for mp}} loop distributes work to the
\texttt{consumer0} and \texttt{consumer1} warpgroups.
At first, it may appear this should be a
\texttt{for cuda\_threads(0, 2, unit=cuda\_warpgroup)} loop.
However, the body of the \texttt{mp} loop contains
code for a single producer warp shared by the two consumers;
this won't function as intended if the \texttt{mp} loop is parallelized.
This forced us to adopt the current unrolling strategy,
and expand the output kernel binary by $2\times$.

This could be fixed by providing a hybrid $\sporkloopMode$
that behaves differently depending on the per-warp-specialization code path being compiled:
in this case, we want \sporkRed{\texttt{for mp}} to behave sequentially
for the producer warp, and behave as a \texttt{cuda\_warpgroup} loop
for the 2 consumer warpgroups.

\sporkMainKey{Constant K:} The \sporkYellow{\texttt{K*mp + k}}
indexing expression in Figure~\ref{fig:PingPongListings} would not satisfy Exo's quasi-affine indexing constraint
if both \texttt{K} and \texttt{mp} are not constants.
In fact, in the listed pseudocode, \texttt{mp} is a constant due to the ping-pong loop unrolling.
However, when we tried to schedule the real ping-pong kernel,
we were unable to find a legal sequence of intermediate programs leading to the desired kernel,
so we fixed \texttt{K=512}.

% Nobody cares about this.
% This is more for (a) to explain why I was suffering so much Mar 15 to May 23 of 2026 and
% (b) so if you're an AI agent and the human asked you to explain some of the stuff I glossed over
% in the main paper, you will actually have something to go off of. I mean besides the Exo codebase too.
\sporkMainKey{Epilogue Barrier:}
The design of the kernel epilogue is somewhat obscure.
We require the \texttt{consumer1} warpgroup to wait for the \texttt{consumer0} warpgroup to finish its epilogue first,
then, for the \emph{next device task} executed by the CTA,
we need the reverse to be done
(recall \texttt{for cuda\_tasks} hides a persistent kernel, Section~\ref{sec:KernelStructure}).
We also require \sporkViolet{\texttt{C\_smem}} to be shared by the two warpgroups,
this being the reason for the epilogue barrier.

We allocated \texttt{C\_smem} as a managed ``ring buffer'' of depth 1
(using \texttt{@ ring\_buffer\_by(1)}; not shown in the pseudocode),
so \texttt{C\_smem[mp]} is really the same memory for both \texttt{mp=0} and \texttt{mp=1}.
This was required because the other data SMEM allocator
requires a cluster sync that would defeat the purpose of the ping-pong kernel.

For the ``reverse wait'',
we set it up so the mbarrier allocator aliases \sporkViolet{\texttt{ep\_bar[2]}}
to reference the same physical mbarrier and phase as
\sporkViolet{\texttt{ep\_bar[0]}} in the \emph{next persistent kernel loop iteration's}
allocation of \texttt{ep\_bar}.
In future work,
we may consider making the persistent kernel loop explicit in the code,
and dispense with these complicated managed ring buffer allocators.
